\section{Conclusion}

We presented a dual-encoder segmentation architecture that fuses RGB photometric features with Canny-derived geometric edge features through multi-scale cross-attention. Three design choices distinguish the proposed model from a naive bottleneck fusion baseline: a learnable scalar gate that allows the model to down-weight noisy edge contributions during training, a second cross-attention module at encoder stage 4 for earlier structural injection, and edge encoder skip connections that propagate boundary cues through the full decoder path.

On the Cataract-Seg dataset, the proposed model achieves an IoU of 0.9530, matching the best-performing baseline (Early Fusion, 0.9532). The results demonstrate that principled multi-scale cross-attention fusion is competitive with simpler fusion strategies, and that the gated residual formulation prevents performance degradation when the auxiliary modality is noisy — a practical advantage in clinical settings where edge map quality may vary with imaging conditions.

The primary limitation of this evaluation is dataset size: with 30 test images, the differences among the top three models are within statistical noise. Future work should validate the proposed architecture on larger ophthalmic datasets and explore learnable edge detectors (e.g., Sobel filters with trainable weights) as a replacement for fixed Canny filtering, which may allow the edge stream to capture more task-relevant structural features. Extension to multi-class segmentation of individual anterior segment structures (cornea, pupil, lens separately) and to intraoperative video frames~\cite{grammatikopoulou2021cadis} are also natural next steps.
